\chapter{The hyperbolic converses: positive, mixed, and exact-profile}
\label{chap:hyperbolic_converse}

% Chapter for the three top-level H² wrapper files:
%   Gluck/Hyperbolic/Converse.lean — positive (escape-velocity) wrappers:
%     RealizesHyperbolicCurvature, HyperbolicFourVertex, hyperbolic_gluck_converse
%     (the ε = −1 instance of SpaceForm.gluck_converse);
%   Gluck/Hyperbolic/Mixed.lean — mixed-sign wrappers:
%     MixedHyperbolicFourVertex, hyperbolic_dahlberg_converse_reparam (up to reparam),
%     hyperbolic_dahlberg_converse (re-exports of the Hyperbolic.* results);
%   Gluck/Hyperbolic/Exact.lean — degree-one reparam removal:
%     injOn_Ico_shift_of_periodic, realizes_of_reparam_degree_one,
%     Hyperbolic.dahlberg_converse (closes the AL-6 gap).
% Absorbs the prose of the retired chapter Gluck_HyperbolicMixed.tex; its two
% content labels def:mixed_hyperbolic_four_vertex and
% thm:hyperbolic_mixed_converse move here verbatim (label + pin preserved).
% SOURCE background: references/spaceform_notes.md ("H^2 vs S^2", the metric
% denominator (1 - |z|^2) > 0 unlocking negative-kappa arcs).

\section{Overview}

This chapter states the three H\textsuperscript{2} $(K=-1)$ four-vertex
converses in their user-facing form, in the Poincar\'e-disk model
$\{|z|<1\}\subseteq\mathbb C$ with metric $g=4/(1-|z|^2)^2\,|dz|^2$. All three
are thin wrappers: the mathematics lives in the space-form engine
(\cref{thm:spaceform_converse_pos}) and the H\textsuperscript{2} arc-length /
fork-A family chain culminating in
\cref{thm:hyperbolic_mixed_converse_reparam_deg1}. The chapter's one genuinely
new proof is the degree-one reparametrization-removal lemma
\cref{thm:realizes_of_reparam_degree_one}, which upgrades the up-to-reparam
mixed converse to the \emph{exact-profile} capstone
\cref{thm:hyperbolic_mixed_converse_exact}.

\begin{itemize}
\item \emph{Positive (escape-velocity) converse},
  \cref{thm:hyperbolic_converse_pos}: a continuous, $2\pi$-periodic $\kappa$
  with the value-separated four-vertex extrema and the escape-velocity bound
  $\kappa>1$ \emph{everywhere} is realized by a simple closed curve. The bound
  is the H\textsuperscript{2} threshold: circles have geodesic curvature $>1$,
  horocycles exactly $1$, and curves with $\kappa\le1$ can run off to infinity,
  so $\kappa>1$ is the natural positivity hypothesis. This is the $K=-1$
  instance of the $K$-generic space-form positive converse; nothing
  hyperbolic-specific remains to prove.
\item \emph{Mixed (Dahlberg) converse},
  \cref{thm:hyperbolic_mixed_converse}: $\kappa$ may dip genuinely negative —
  in fact \emph{arbitrarily} negative (\cref{sec:hyperbolic_negative_scope}) —
  provided the extrema alternate in the value-separated four-vertex pattern
  with the two maxima above an escape level $c>1$. The conclusion is stated up
  to an orientation-preserving $C^1$ reparametrization $\Psi$.
\item \emph{Exact-profile converse} — the headline of the chapter —
  \cref{thm:hyperbolic_mixed_converse_exact}: the same hypothesis realizes
  $\kappa$ \emph{exactly},
  \[
    \exists\,\gamma,\quad \texttt{IsSimpleClosed }\gamma\ \wedge\
    \texttt{Realizes }(-1)\ \gamma\ \kappa,
  \]
  with \emph{no} reparametrization. The up-to-reparam form
  \cref{thm:hyperbolic_mixed_converse} is kept as a weakening corollary.
\end{itemize}

\paragraph{Exact versus up-to-reparam.}
The distinction matters. ``Up to reparametrization'' realizes $\kappa\circ\Psi$
for \emph{some} orientation-preserving $C^1$ change of parameter $\Psi$ — the
curve traverses the same curvature values in the same cyclic order, but the
prescribed function itself is only recovered after an unknown distortion of
the parameter. ``Exact'' realizes the given $\kappa$ on the nose, matching the
statements of the Euclidean \cref{thm:dahlberg_converse} and the spherical
\cref{thm:spherical_converse}. The reason exact is reachable on
H\textsuperscript{2} is that the fork-A reparametrization is not an arbitrary
increasing map: the family chain produces
$\Psi=(h_1\circ\mathrm{nodeMap})\circ\chi$, and the node map conjugates the
co-constructed arc-length period $\Lambda$ to $2\pi$
(\cref{lem:node_map_add_period}), so the composite is a \emph{degree-one
circle map}: $\Psi(t+2\pi)=\Psi(t)+2\pi$. This is recorded in
\cref{thm:hyperbolic_mixed_converse_reparam_deg1} (the up-to-reparam converse
with the degree-one clause attached), and
\cref{thm:realizes_of_reparam_degree_one} shows that a degree-one,
orientation-preserving $C^1$ reparametrization can always be removed: the
global $C^1$ inverse $\Psi^{-1}$ exists, and $w=\gamma\circ\Psi^{-1}$ is a simple
closed curve realizing $\kappa$ exactly.

\paragraph{Historical honesty.}
The up-to-reparam form of \cref{thm:hyperbolic_mixed_converse} was an artifact
of the \emph{abstract} window analysis, not of the geometry.
H\textsuperscript{2} has no metric rescaling, so the arc-length period
$\Lambda$ is co-constructed rather than normalized to $2\pi$, and the abstract
converse could only assert an unstructured $C^1$ reparametrization. But the
\emph{concrete} node map of the fork-A construction conjugates $\Lambda$ to
$2\pi$ exactly (\cref{lem:node_map_add_period}); once that degree-one
structure of $\Psi$ was surfaced through
\cref{thm:hyperbolic_mixed_converse_reparam_deg1}, the reparametrization could
be removed outright. This closed the AL-6 gap.

\paragraph{Relation to the space-form wrappers.}
The $K$-generic packaging of these results lives in
\cref{chap:spaceform_converse}: in particular the space-form-level hyperbolic
Dahlberg converse \cref{thm:hyperbolic_dahlberg_converse} is derived
\emph{from} the exact capstone \cref{thm:hyperbolic_mixed_converse_exact} of
this chapter, not the other way around — the flow of implication is
family chain $\to$ exact profile $\to$ space-form restatement.

\section{Scope: arbitrarily negative minima}
\label{sec:hyperbolic_negative_scope}

The minima are \emph{arbitrarily negative} --- there is no lower bound on
$\kappa$. The hypothesis constrains only the maxima, through the value-separation
$\max(1,\kappa(q_1),\kappa(q_2))<\min(\kappa(p_1),\kappa(p_2))$ and a window level
$c>1$; the minima $\kappa(q_1),\kappa(q_2)$ may be as negative as one likes. An
earlier version carried a confinement floor
$-\bigl(\operatorname{centeredRadius}(-1)\,c\bigr)=-\bigl(c-\sqrt{c^2-1}\bigr)$,
but it was vestigial: the fork-A route reaches the curve through convex clean
levels $1<a<b$ and the $L^1$-closeness of $\kappa\circ h_1$ to a convex reference
bicircle, which absorbs dips of any depth (Dahlberg's $L^1$ squeeze). Removing it
gives the \emph{full genuinely-negative} converse, strictly larger than the
everywhere-escape positive converse (which it subsumes via
\texttt{MixedSignHyperbolicFourVertex.of\_escape\_positive}). Unlike the
spherical case, the H\textsuperscript{2} metric denominator
$1-\lVert z\rVert^2>0$ is a confinement condition, not an admissibility one,
so concave / negative-$\kappa$ arcs are permitted by the reconstruction flow
without a positive-speed obstruction.

\section{Declarations}

\begin{definition}\leanok
[hyperbolic realization]
  \label{def:realizes_hyperbolic_curvature}
  \lean{Gluck.RealizesHyperbolicCurvature}
  \uses{def:spaceform_realizes}
  A curve $\gamma:\mathbb R\to\mathbb C$ \emph{realizes the hyperbolic curvature
  function} $\kappa$ if $\texttt{SpaceForm.Realizes}\,(-1)\,\gamma\,\kappa$: against
  the tangent-angle gauge $\varphi$, the defining speed relation is
  \[
    \frac{1-\lVert \gamma\rVert^2}{2}\,\varphi'
      \;=\;\bigl(\kappa+\langle \gamma,\,ie^{i\varphi}\rangle_{\mathbb R}\bigr)
      \,\lVert \gamma'\rVert
  \]
  (note the $+$ inner-product term, versus the $-$ of the spherical
  \cref{def:realizes_spherical_curvature}, and the metric factor
  $1-\lVert z\rVert^2$). Definitionally the $K=-1$ instance of the
  $K$-generic predicate; the accompanying lemma
  \texttt{realizesHyperbolicCurvature\_iff\_realizes\_neg\_one} records the
  (reflexive) unfolding.
\end{definition}

\begin{definition}\leanok
[hyperbolic four-vertex hypothesis]
  \label{def:hyperbolic_four_vertex}
  \lean{Gluck.HyperbolicFourVertex}
  \uses{def:is_curvature_function, def:four_vertex_condition}
  The $K=-1$ instance of \texttt{SpaceForm.SpaceFormFourVertex}: a
  continuous, $2\pi$-periodic, strictly positive $\kappa$ with the
  value-separated four-vertex extrema (\cref{def:four_vertex_condition}) and
  the escape-velocity bound $\kappa>1$ everywhere. Since the ambient sign
  $K=-1$ is negative, the conditional escape guard of the generic
  predicate fires unconditionally, and the unfolding lemma
  \texttt{hyperbolicFourVertex\_iff} states the hypothesis as the plain
  conjunction: $\kappa$ is a curvature function, satisfies the four-vertex
  condition, and $1<\kappa(\theta)$ for all $\theta$.
\end{definition}

\begin{theorem}\leanok
[positive hyperbolic converse]
  \label{thm:hyperbolic_converse_pos}
  \lean{Gluck.hyperbolic_gluck_converse}
  \uses{def:hyperbolic_four_vertex, def:realizes_hyperbolic_curvature,
        def:is_simple_closed, thm:spaceform_converse_pos}
  A hyperbolic four-vertex curvature function
  (\cref{def:hyperbolic_four_vertex}) is realized by a simple closed curve in
  the Poincar\'e disk:
  \[
    \texttt{HyperbolicFourVertex }\kappa\;\Longrightarrow\;
    \exists\,\gamma,\ \texttt{IsSimpleClosed }\gamma\ \wedge\
    \texttt{RealizesHyperbolicCurvature }\gamma\,\kappa .
  \]
  Genuinely new mathematics: there is no published Gluck--Dahlberg converse
  for prescribed cyclic geodesic curvature on H\textsuperscript{2}.
\end{theorem}

\begin{proof}
  \uses{thm:spaceform_converse_pos}
  \leanok
  One line: the $K=-1$ case of the space-form positive converse,
  $\texttt{SpaceForm.gluck\_converse}\ (\texttt{Or.inr rfl})$. All the
  work — flow, winding, reconstruction, simplicity — happens
  $K$-generically in \cref{chap:spaceform_converse}.
\end{proof}

\begin{definition}[Mixed-sign hyperbolic four-vertex hypothesis]
  \label{def:mixed_hyperbolic_four_vertex}
  \lean{Gluck.MixedHyperbolicFourVertex}
  \leanok
  \uses{def:mixed_sign_hyperbolic_four_vertex, def:four_vertex_condition}
  A continuous, $2\pi$-periodic $\kappa$ that is either constant escape
  ($\exists c,\ 1<c\ \wedge\ \kappa\equiv c$) or has the value-separated
  alternating four-vertex extrema with both maxima above an escape level $c>1$.
  No lower bound on the minima. Re-export of the working predicate
  \cref{def:mixed_sign_hyperbolic_four_vertex}
  (\texttt{Hyperbolic.MixedSignHyperbolicFourVertex}), mirroring how
  \cref{def:hyperbolic_four_vertex} fronts the space-form predicate; the
  unfolding lemma \texttt{mixedHyperbolicFourVertex\_iff\_mixedSign} is
  reflexive.
\end{definition}

\begin{theorem}[Hyperbolic mixed converse, up to reparametrization]
  \label{thm:hyperbolic_mixed_converse}
  \lean{Gluck.hyperbolic_dahlberg_converse_reparam}
  \leanok
  \uses{def:mixed_hyperbolic_four_vertex, def:realizes_hyperbolic_curvature,
        def:is_simple_closed}
  A genuinely-negative mixed-sign hyperbolic four-vertex curvature function is
  realized, up to an orientation-preserving $C^1$ reparametrization $\Psi$
  ($0<\Psi'$), as the hyperbolic geodesic curvature of a simple closed curve in
  the Poincar\'e disk:
  \[
    \exists\,\gamma\,\Psi,\quad C^1\,\Psi\ \wedge\ 0<\Psi'\ \wedge\
    \texttt{IsSimpleClosed }\gamma\ \wedge\
    \texttt{RealizesHyperbolicCurvature }\gamma\,(\kappa\circ\Psi).
  \]
  As with the arc-length family results, the conclusion is stated up to
  reparametrization because H\textsuperscript{2} has no metric rescaling: the
  period is co-constructed rather than normalized. Retained as a weakening
  corollary of the exact capstone
  \cref{thm:hyperbolic_mixed_converse_exact}; see the historical note in the
  overview.
\end{theorem}

\begin{proof}
  \uses{thm:hyperbolic_mixed_converse_reparam_deg1}
  \leanok
  Re-export of the family-chain assembly
  (\texttt{Hyperbolic.dahlberg\_converse\_reparam},
  \cref{chap:hyperbolic_family}): symbolic $(a,b)$-family anchor, five-leg
  Gr\"onwall transport, the asymmetric $3$-degree-of-freedom
  Poincar\'e--Miranda closing, the radial-monotonicity simplicity transport,
  and the window-bridge reindexing from the flow horizon $2L$ to the profile
  period $\Lambda=\mathrm{nodePeriod}\le 2L$. A key structural point is that
  the closing/simplicity constants $C_1,\varepsilon_0,\mu$ are quantified
  \emph{uniformly over the reparametrization} $h_1$, so the assembled
  tolerance $\varepsilon:=\min(\varepsilon_0,\mu/C_1)$ can be fixed
  \emph{before} the $L^1$ reparametrization is built at tolerance
  $\varepsilon$, breaking the reparam/$\varepsilon$ fixed-point circularity.
  Equivalently: weaken \cref{thm:hyperbolic_mixed_converse_reparam_deg1} by
  forgetting its degree-one clause.
\end{proof}

\begin{lemma}\leanok
[injectivity on shifted fundamental intervals]
  \label{lem:injon_ico_shift_of_periodic}
  \lean{Gluck.Hyperbolic.injOn_Ico_shift_of_periodic}
  Let $f:\mathbb R\to\mathbb C$ be $T$-periodic ($T>0$) and injective on the
  fundamental interval $[0,T)$. Then $f$ is injective on every shifted
  fundamental interval $[c,c+T)$.
\end{lemma}

\begin{proof}

\leanok
  Reduce each point into $[0,T)$ by subtracting $\lfloor x/T\rfloor\,T$
  (periodicity preserves the value; membership is
  $x-\lfloor x/T\rfloor T=T\{x/T\}\in[0,T)$). If $f(x)=f(y)$ with
  $x,y\in[c,c+T)$, injectivity on $[0,T)$ forces the reduced points to agree,
  so $x-y=kT$ for the integer $k=\lfloor x/T\rfloor-\lfloor y/T\rfloor$; but
  $|x-y|<T$ within a half-open interval of length $T$, so $k=0$ and $x=y$.
\end{proof}

\begin{theorem}\leanok
[reparam removal for degree-one reparametrizations]
  \label{thm:realizes_of_reparam_degree_one}
  \lean{Gluck.Hyperbolic.realizes_of_reparam_degree_one}
  \uses{def:spaceform_realizes, def:is_simple_closed,
        lem:injon_ico_shift_of_periodic}
  Let $K\in\mathbb R$ and suppose a simple closed curve $\gamma$ realizes
  $\kappa\circ\Psi$ at curvature $K$, where $\Psi$ is $C^1$,
  orientation-preserving ($0<\Psi'$ everywhere), and a \emph{degree-one circle
  map}: $\Psi(t+2\pi)=\Psi(t)+2\pi$. Then $\kappa$ is realized \emph{exactly}
  by the simple closed curve $w=\gamma\circ\Psi^{-1}$:
  \[
    \exists\,w,\quad \texttt{IsSimpleClosed }w\ \wedge\
    \texttt{Realizes }K\ w\ \kappa .
  \]
  Stated $K$-generically: nothing in the removal is hyperbolic.
\end{theorem}

\begin{proof}
  \uses{lem:injon_ico_shift_of_periodic}
  \leanok
  \emph{Global inverse.} $\Psi$ is strictly monotone
  (\texttt{strictMono\_of\_deriv\_pos}), and the degree-one identity iterates
  to $\Psi(x+2\pi n)=\Psi(x)+2\pi n$, forcing $\Psi\to\pm\infty$ at
  $\pm\infty$; with continuity this gives surjectivity, so $\Psi$ packages as
  an order isomorphism / self-homeomorphism $e$ of $\mathbb R$ with inverse
  $\Phi=\Psi^{-1}$.

  \emph{$\Phi$ is a $C^1$, orientation-preserving degree-one map.} The global
  inverse-function theorem for homeomorphisms with nonvanishing derivative —
  mathlib's \texttt{Homeomorph.contDiff\_symm\_deriv} applied to $e$ with the
  witnesses $\Psi'\ne0$ — gives $\Phi\in C^1$;
  \texttt{HasDerivAt.of\_local\_left\_inverse} computes
  $\Phi'(t)=\bigl(\Psi'(\Phi t)\bigr)^{-1}>0$; and applying the injective
  $\Psi$ to both sides of the degree-one identity transfers it to $\Phi$:
  $\Phi(t+2\pi)=\Phi(t)+2\pi$.

  \emph{$w=\gamma\circ\Phi$ is simple closed.} Closedness: $w(t+2\pi)
  =\gamma(\Phi(t)+2\pi)=\gamma(\Phi(t))=w(t)$ by the degree-one identity for $\Phi$ and
  $2\pi$-periodicity of $\gamma$. Injectivity on $[0,2\pi)$: $\Phi$ maps
  $[0,2\pi)$ strictly monotonically into the \emph{shifted} fundamental
  interval $[\Phi(0),\Phi(0)+2\pi)$, on which $\gamma$ is injective by
  \cref{lem:injon_ico_shift_of_periodic} — this injectivity transfer across a
  shifted period is exactly why the helper lemma exists.

  \emph{$w$ realizes $\kappa$.} The reparametrization-transport lemma
  \texttt{spaceFormRealizes\_comp} (\cref{chap:spaceform_converse}) applied to
  the realization of $\kappa\circ\Psi$ along the $C^1$, $\Phi'>0$ map $\Phi$
  gives $\texttt{Realizes }K\,(\gamma\circ\Phi)\,((\kappa\circ\Psi)\circ\Phi)$,
  and $(\kappa\circ\Psi)\circ\Phi=\kappa$ by the right-inverse identity
  $\Psi(\Phi t)=t$.
\end{proof}

\begin{theorem}\leanok
[exact-profile hyperbolic mixed converse]
  \label{thm:hyperbolic_mixed_converse_exact}
  \lean{Gluck.hyperbolic_dahlberg_converse}
  \uses{def:mixed_hyperbolic_four_vertex, def:realizes_hyperbolic_curvature,
        def:is_simple_closed, thm:hyperbolic_mixed_converse_reparam_deg1,
        thm:realizes_of_reparam_degree_one}
  A genuinely-negative mixed-sign hyperbolic four-vertex curvature function
  (\cref{def:mixed_hyperbolic_four_vertex}) is realized \emph{exactly} — with
  no reparametrization — as the hyperbolic geodesic curvature of a simple
  closed curve in the Poincar\'e disk:
  \[
    \texttt{MixedHyperbolicFourVertex }\kappa\;\Longrightarrow\;
    \exists\,\gamma,\ \texttt{IsSimpleClosed }\gamma\ \wedge\
    \texttt{Realizes }(-1)\ \gamma\ \kappa .
  \]
  The H\textsuperscript{2} capstone: it subsumes both
  \cref{thm:hyperbolic_converse_pos} (whose hypothesis implies the mixed one
  via \texttt{MixedSignHyperbolicFourVertex.of\_escape\_positive}) and the
  up-to-reparam \cref{thm:hyperbolic_mixed_converse}, and it is the statement
  from which the space-form-level \cref{thm:hyperbolic_dahlberg_converse} is
  derived in \cref{chap:spaceform_converse}. It matches the exact-profile
  shape of the Euclidean \cref{thm:dahlberg_converse} and the spherical
  \cref{thm:spherical_converse}.
\end{theorem}

\begin{proof}
  \uses{thm:hyperbolic_mixed_converse_reparam_deg1,
        thm:realizes_of_reparam_degree_one}
  \leanok
  Two steps. \cref{thm:hyperbolic_mixed_converse_reparam_deg1}
  (\cref{chap:hyperbolic_family}) supplies a simple closed $\gamma$ realizing
  $\kappa\circ\Psi$ where the fork-A reparametrization
  $\Psi=(h_1\circ\mathrm{nodeMap})\circ\chi$ is $C^1$ with $0<\Psi'$ \emph{and}
  degree-one, $\Psi(t+2\pi)=\Psi(t)+2\pi$ — the degree-one clause is genuine
  because the node map conjugates the co-constructed period $\Lambda$ to
  $2\pi$ (\cref{lem:node_map_add_period}). Then
  \cref{thm:realizes_of_reparam_degree_one} removes $\Psi$ via
  $w=\gamma\circ\Psi^{-1}$. The user-facing \texttt{Gluck.hyperbolic\_dahlberg\_converse}
  is the verbatim re-export of this composite through
  \cref{def:mixed_hyperbolic_four_vertex}.
\end{proof}
